% !TeX root = ../../Book1.tex
% 纲要标题：## 第11章　整除、主理想整环与唯一分解
% 纲要位置：计划纲要.md，第 414 行。

\chapter{\texorpdfstring{\mbox{整除、主理想整环}}{整除、主理想整环}\texorpdfstring{\mbox{与唯一分解}}{与唯一分解}}
\label{b1:ch:11}
\BookChapterNumber{b1:ch:11}

\begin{chapterabstract}
整数的素因数分解在一般整环中需要重新检验：不可约元未必为素元，最大公因子也未必给出 Bézout 等式。本章比较 Euclid 整环、主理想整环与唯一分解整环，说明它们的蕴含关系、反例和可执行的整除计算。
\end{chapterabstract}

\begin{learninggoals}
\begin{itemize}
\item 区分不可约元与素元，区分最大公因子与 Bézout 关系。
\item 在整数、域上的一元多项式和 Gauss 整数中执行 Euclid 算法及回代。
\item 证明 Euclid 整环是主理想整环、主理想整环是唯一分解整环，并指出分解存在性与唯一性各自的依据。
\item 用具体理想与范数计算检验主性或唯一分解性的失败，理解理想分解问题的来源。
\end{itemize}
\end{learninggoals}

在整数中，\(2\) 与 \(3\) 的最大公因子是 \(1\)，而且能找到整数 \(u,v\) 使 \(2u+3v=1\)。把同样的话移到 \(\mathbb Z[x]\)，\(2\) 与 \(x\) 仍没有非单位公因子，却不可能有 \(2u(x)+xv(x)=1\)：令 \(x=0\) 就会得到偶数等于 \(1\)。最大公因子、Bézout 等式和理想由一个元素生成，在一般整环中不再是同一回事。本章从整除与 Euclid 算法出发，分辨这些性质之间真正成立的推论，再考察元素何时能唯一分解。

\ProseParagraphBreak

本章使用\chapref{b1:ch:10}建立的整环、单位、理想、商环及多项式次数性质；Gauss 整数和非 Euclid 的例子还用到复数的共轭与模。非 Euclid 主理想整环的例子用于辨析条件强弱。\(\mathbb Z[x]\) 的非主性在本章直接证明，唯一分解性则明确留待下一章的多项式定理。

\ProseParagraphBreak
整除、主理想与唯一分解，可结合 Dummit 与 Foote 的《Abstract Algebra》\cite[Sections 8.1--8.3]{DummitFoote2003}学习。整数的 Euclid 算法、扩展算法与模逆可读 Shoup 的《A Computational Introduction to Number Theory and Algebra》\cite[Sections 4.1--4.3]{Shoup2008NumberTheoryAlgebra}；域上一元多项式的 Euclid 算法、模逆及中国剩余计算见同书\cite[Sections 17.3--17.4]{Shoup2008NumberTheoryAlgebra}。唯一分解整环、Euclid 整环与主理想整环的结构性讨论见该书\cite[Section 16.9.1]{Shoup2008NumberTheoryAlgebra}。

% !TeX root = ../../Book1.tex
% 纲要标题：### 11.1 整除与因子
% 纲要位置：计划纲要.md，第 416 行。

\section{整除与因子}
\label{b1:sec:1101}

% 纲要内容（仅作写作计划，尚非教材正文）：
%
% * 相伴、不可约元与素元
% * 最大公因子与 Bézout 关系
% * 元素分解与理想语言的区别
%

整数的素因数分解同时包含两个断言：每个非零非单位都能分解，所得不可约因子又在适当意义下唯一。在一般整环中，这两个断言必须分别考察。本节先澄清整除、不可约与素性之间的关系；本章的环均为交换整环，除非另有说明。

\subsection{相伴、不可约元与素元}
\label{b1:subsec:1101:01}

\begin{definition}\label{b1:def:divisibility-associates}
设 \(R\) 是整环，\(a,b\in R\)。若存在 \(c\in R\) 使 \(b=ac\)，则称 \(a\) 整除 \(b\)，记为 \(a\mid b\)；这一关系称为\term[zhengchu]{整除关系}[divisibility]。若存在单位 \(u\in R^\times\) 使 \(a=ub\)，则称 \(a,b\) \term[xiangban]{相伴}[associate]。
\end{definition}
\symidx{\(a\mid b\)：\(a\) 整除 \(b\)}

每个元素都整除 \(0\)，而 \(0\) 只整除 \(0\)。单位整除所有元素。相伴是等价关系：单位 \(1\)、\(u^{-1}\) 和单位的乘积分别给出自反性、对称性和传递性。在 \(\mathbb Z\) 中，相伴就是相差一个符号；在域中，任意两个非零元素都相伴。

\begin{proposition}\label{b1:prop:associate-ideals}
设 \(R\) 为整环，\(a,b\in R\)。则 \(a\mid b\) 当且仅当 \((b)\subseteq(a)\)。此外，以下三条件等价：\(a,b\) 相伴；\(a\mid b\) 且 \(b\mid a\)；\((a)=(b)\)。
\end{proposition}
\begin{proof}
若 \(b=ac\)，则 \(br=a(cr)\)，故 \((b)\subseteq(a)\)。反之，该包含使 \(b\in(a)\)，于是 \(b=ac\)。相伴当然给出相互整除。若 \(a=bc\)、\(b=ad\) 且 \(a\ne0\)，则 \(a=acd\)，整环中的消去律给出 \(cd=1\)，故 \(c\) 为单位。若 \(a=0\)，相互整除迫使 \(b=0\)，两者也相伴。主理想相等与相互整除的等价性来自第一句话。
\end{proof}

\begin{definition}\label{b1:def:irreducible-prime}
设 \(p\) 是整环 \(R\) 中的非零非单位。
\begin{enumerate}
\item 若对每个满足 \(p=ab\) 的 \(a,b\in R\)，二者中至少一个为单位，则称 \(p\) 为\term[bukeyueyuan]{不可约元}[irreducible element]。
\item 若对任意 \(a,b\in R\)，\(p\mid ab\) 都蕴含 \(p\mid a\) 或 \(p\mid b\)，则称 \(p\) 为\term[suyuan]{素元}[prime element]。
\end{enumerate}
\end{definition}

不可约性限制 \(p\) 自身的分解，素性涉及 \(p\) 对任意乘积的整除。两者并非同一个定义。单位被排除，是因为在分解中插入或删去单位不应改变因子的实质。

\begin{proposition}\label{b1:prop:prime-implies-irreducible}
设 \(R\) 为整环，\(p\in R\) 为非零非单位。若 \(p\) 是素元，则 \(p\) 是不可约元；而且 \(p\) 是素元，当且仅当主理想 \((p)\) 是 \(R\) 的素理想。
\end{proposition}
\begin{proof}
设 \(p=ab\)。素性给出 \(p\mid a\) 或 \(p\mid b\)。若 \(a=pc\)，则 \(p=pcb\)，由 \(p\ne0\) 消去得到 \(cb=1\)，所以 \(b\) 为单位；另一种情形交换 \(a,b\) 即得。又因 \(p\) 非单位，\((p)\ne R\)；条件 \(ab\in(p)\) 蕴含 \(a\in(p)\) 或 \(b\in(p)\)，逐项翻译后正是素元的定义。
\end{proof}

后面将证明，在主理想整环和唯一分解整环中，不可约元也都是素元；同时会看到 \(\mathbb Z[\sqrt{-5}]\) 中的不可约元 \(2\) 并不是素元。此处不能把整数中的直觉当作一般整环的结论。

\subsection{最大公因子与 Bézout 关系}
\label{b1:subsec:1101:02}

\begin{definition}\label{b1:def:gcd}
设 \(a,b\) 为整环 \(R\) 的两个元素。如果 \(d\in R\) 满足 \(d\mid a\)、\(d\mid b\)，并且每个同时整除 \(a,b\) 的元素 \(c\in R\) 都整除 \(d\)，就称 \(d\) 为 \(a,b\) 的\term[zuidagongyinzi]{最大公因子}[greatest common divisor]。如果 \(1\) 是 \(a,b\) 的最大公因子，就称 \(a,b\) \term[husu-yuansu]{互素}[relatively prime]。
\end{definition}

这里的“最大”依整除关系而言，与实数大小无关。最大公因子若存在，由\cref{b1:prop:associate-ideals}可知它在相伴意义下唯一；未经额外约定，\(\gcd(a,b)\) 表示任选的一个最大公因子。在整数中通常取非负者，在域上的多项式环中通常取首一者；约定 \(\gcd(0,0)=0\)。有限多个元素的最大公因子同样按“公共因子都整除它”定义。
\symidx{\(\gcd(a,b)\)：最大公因子（相伴意义下）}

\begin{proposition}\label{b1:prop:bezout-gcd}
设 \(R\) 为整环，\(a,b,d\in R\)。若 \(d\mid a\)、\(d\mid b\)，且存在 \(u,v\in R\) 使
\begin{equation}\label{b1:eq:bezout-ring}
d=ua+vb,
\end{equation}
则 \(d\) 为 \(a,b\) 的最大公因子。特别地，\((a,b)=R\) 蕴含 \(a,b\) 互素。
\end{proposition}
\begin{proof}
任一公因子 \(c\) 整除 \(a,b\)，于是也整除线性组合 \(ua+vb=d\)。这与已给出的 \(d\mid a,b\) 正好构成最大公因子的两项要求。若 \((a,b)=R\)，则 \(1=ua+vb\)；又 \(1\) 整除一切元素，所以取 \(d=1\) 即得。
\end{proof}

设 \(R\) 为整环，\(a,b\in R\)，并设 \(d\) 是 \(a,b\) 的最大公因子。若存在 \(u,v\in R\) 使 \(d=ua+vb\)，则称这个等式为 \(a,b\) 的\term[bezout-guanxi]{Bézout 关系}[Bézout identity]。它比“\(d\) 是最大公因子”多提供了一个线性组合表达式。在 \(\mathbb Z\) 中，\(\gcd(30,18)=6\)，而 \(6=2\cdot18-30\) 给出相应关系。一般整环中，即使最大公因子存在，也未必能如此表示；下面的例子正说明这一点。

\subsection{元素分解与理想}
\label{b1:subsec:1101:03}

整除关系在主理想之间反向排列，因此“共同整除”与“共同生成”不可混淆。对两个非零元素 \(a,b\)，若 \((a,b)=(d)\)，则 \(d\) 整除 \(a,b\)，且 \(d\) 属于它们生成的理想，\cref{b1:prop:bezout-gcd}遂给出最大公因子和 Bézout 关系。反过来，知道最大公因子并不保证这个二生成理想是主理想。

\begin{proposition}\label{b1:prop:integer-poly-nonprincipal}
在 \(\mathbb Z[x]\) 中，\(2,x\) 的最大公因子为 \(1\)，但 \((2,x)\) 是非主理想，特别地不等于全环。
\end{proposition}
\begin{proof}
若非零多项式 \(c\) 同时整除 \(2,x\)，由\cref{b1:prop:polynomial-degree}的次数相加公式，\(c\mid2\) 迫使 \(c\) 是常数，且为 \(\pm1\) 或 \(\pm2\)。后两者不整除 \(x\)，因为 \(x\) 的一次项系数为 \(1\)。故 \(\gcd(2,x)=1\)。然而若有 \(1=2f(x)+x\cdot g(x)\)，比较常数项就得到 \(1=2f(0)\)，不可能。因此 \((2,x)\ne\mathbb Z[x]\)。若 \((2,x)=(h)\)，则 \(h\) 是 \(2,x\) 的公因子，因而为单位；这又会使该理想等于全环，矛盾。
\end{proof}

互素并不蕴含生成单位理想；一个理想仅用两个元素生成，也可能无法缩减成一个生成元。主理想整环要求每个理想都是主理想，不论是否预先知道它有限生成，因此排除了这种情形。仅要求二生成理想都是主理想，只能进一步推出所有有限生成理想都是主理想：对生成元的个数作归纳，每次将已有的一个生成元与下一个合并即可；这一条件还不足以保证所有理想都是主理想。\cref{b1:ex:11-factorization-nonexistence}将给出一个所有有限生成理想都为主、却不是主理想整环的例子。

\ProseParagraphBreak
不可约性也可以用主理想说明。设 \(R\) 为整环，\(p\in R\) 为非零非单位。由整除与主理想包含的反向对应可知，\(p\) 不可约，当且仅当 \((p)\) 在 \(R\) 的真主理想中为极大元：若 \((p)\subseteq(d)\subsetneq R\)，则 \(p=dc\)，其中 \(d\) 非单位；不可约性迫使 \(c\) 为单位，故 \((d)=(p)\)。反过来，若 \(p=ab\) 且 \(a,b\) 都非单位，则 \((p)\subsetneq(a)\subsetneq R\)，与这一极大性矛盾。这里比较的只是主理想；后面在主理想整环中，所有理想都可纳入比较，才得到 \((p)\) 是极大理想。

\ProseParagraphBreak
后文将证明，在 \(R=\mathbb Z[\sqrt{-5}]\) 中，\(2\) 不可约，而非主理想 \(I=(2,1+\sqrt{-5})\) 满足 \((2)\subsetneq I\subsetneq R\)；这与 \((2)\) 在真主理想中极大并不矛盾。

% !TeX root = ../../Book1.tex
% 纲要标题：### 11.2 Euclid 整环
% 纲要位置：计划纲要.md，第 422 行。

\section{Euclid 整环}
\label{b1:sec:1102}

% 纲要内容（仅作写作计划，尚非教材正文）：
%
% * 除法算法与 Euclid 算法
% * 整数环及一元多项式环
% * Gauss 整数的例子
%

整数中求最大公因子的辗转相除法，依靠的是余数不断变小。推广这个方法时，真正需要的不是元素本身有大小，而是给非零余数配一个不能无限下降的非负整数。本节把这一要求写成公理，并在整数、多项式和复平面格点中实际执行除法。

\subsection{除法算法与 Euclid 算法}
\label{b1:subsec:1102:01}

\begin{definition}\label{b1:def:euclidean-domain}
设 \(R\) 是交换整环。若存在函数
\[
\delta:R\setminus\{0\}\longrightarrow\mathbb N
\]
使任意 \(a,b\in R\)、\(b\ne0\) 都可写成 \(a=bq+r\)，其中 \(q,r\in R\)，并且 \(r=0\) 或 \(\delta(r)<\delta(b)\)，则称 \(R\) 为\term[euclid-zhenghuan]{Euclid 整环}[Euclidean domain]，称 \(\delta\) 为\term[euclid-hanshu]{Euclid 函数}[Euclidean function]。
\end{definition}

本书不额外要求 \(\delta(a)\leq\delta(ab)\)（\(a,b\ne0\)）。有些文献把这个乘法单调性并入定义。两种约定对函数的要求不同，但得到的 Euclid 整环类相同。事实上，给定满足上述除法条件的 \(\delta\)，可对每个 \(a\in R\setminus\{0\}\) 定义
\[
\widetilde\delta(a)\coloneqq\min_{c\in R\setminus\{0\}}\delta(ac).
\]
由于 \(R\) 是整环，取最小值的集合是非空的非负整数集合。若 \(a,b\ne0\)，\(ab\) 的每个非零倍数也是 \(a\) 的非零倍数，故
\(\widetilde\delta(a)\leq\widetilde\delta(ab)\)。又设 \(a\in R\)、\(b\in R\setminus\{0\}\)，取 \(c\ne0\) 使 \(\delta(bc)=\widetilde\delta(b)\)。用 \(\delta\) 作带余除法，得到
\[
a=(bc)q+r=b(cq)+r,
\qquad r=0\ \text{或}\ \delta(r)<\delta(bc).
\]
当 \(r\ne0\) 时，在最小值的定义中取乘子 \(1\)，可知 \(\widetilde\delta(r)\leq\delta(r)<\widetilde\delta(b)\)。因此 \(\widetilde\delta\) 同时满足除法条件与乘法单调性\cite[Theorem 2.1]{ConradRemarksEuclideanDomains}。商和余数通常不唯一，Euclid 函数也未必唯一；对带余除法的存在性证明，还须另行判断是否给出了可执行的求商步骤。

\ProseParagraphBreak
设 \(R\) 为带 Euclid 函数 \(\delta\) 的 Euclid 整环，\(a,b\in R\) 不全为零。从 \(a,b\) 出发，不断以上一次非零余数作除数，直到余数为零；每次所得非零余数的 \(\delta\) 值都严格下降。这个过程称为\term[euclid-suanfa]{Euclid 算法}[Euclidean algorithm]，也称辗转相除法。
\termidx[zhanzhuanxiangchufa]{辗转相除法}

\begin{proposition}[Euclid 算法]\label{b1:prop:euclidean-algorithm}
设 \(R\) 为 Euclid 整环，\(a,b\in R\) 不全为零。依次作带余除法，直到余数为零，则最后一个非零余数 \(d\) 是 \(a,b\) 的最大公因子，并可由回代求得 \(u,v\in R\)，使 \(d=ua+vb\)。
\end{proposition}
\begin{proof}
若 \(b=0\)，取 \(d=a\)、\(u=1\)、\(v=0\)。否则置 \(r_0=a,r_1=b\)，并在 \(r_j\ne0\) 时选取
\[
r_{j-1}=q_jr_j+r_{j+1},
\qquad r_{j+1}=0\ \text{或}\ \delta(r_{j+1})<\delta(r_j).
\]
非负整数 \(\delta(r_1),\delta(r_2),\ldots\) 不可能无限严格下降，所以某一步得到 \(r_{m+1}=0\)、\(r_m\ne0\)。等式表明，一个元素同时整除 \(r_{j-1},r_j\)，当且仅当它同时整除 \(r_j,r_{j+1}\)：两个方向分别使用减法和加法。因此每一对相邻余数的公因子集合相同，最终等于 \(r_m,0\) 的公因子集合，故 \(d=r_m\) 为所求最大公因子。

为明确回代，令 \((u_0,v_0)=(1,0)\)、\((u_1,v_1)=(0,1)\)，并递推
\begin{equation}\label{b1:eq:extended-euclid-recurrence}
u_{j+1}=u_{j-1}-q_ju_j,\qquad
v_{j+1}=v_{j-1}-q_jv_j.
\end{equation}
由 \(r_0=u_0a+v_0b\)、\(r_1=u_1a+v_1b\) 开始归纳，除法等式给出每个 \(r_j=u_ja+v_jb\)。取 \(j=m\) 即得所求关系。
\end{proof}

在求余数的同时按\eqref{b1:eq:extended-euclid-recurrence}更新两组系数，便得到\term[kuozhan-euclid-suanfa]{扩展 Euclid 算法}[extended Euclidean algorithm]。\cref{b1:alg:extended-euclid}同时输出最大公因子与 Bézout 系数。Euclid 整环的定义只保证所需的商和余数存在；要实际运行算法，还须能有效完成环运算和符合下降条件的带余除法。

\begin{algorithm}[htbp]
\caption{扩展 Euclid 算法}
\label{b1:alg:extended-euclid}
\begin{algorithmsteps}{配有 Euclid 函数 \(\delta\) 的 Euclid 整环 \(R\)，及不全为零的 \(a,b\in R\)；其中环运算与符合 \(\delta\) 下降条件的带余除法均可有效计算。}{\(a,b\) 的一个最大公因子 \(d\) 及 \(u,v\in R\)，使 \(d=ua+vb\)。}
\State 置 \((r_0,u_0,v_0)=(a,1,0)\)、\((r_1,u_1,v_1)=(b,0,1)\)，并令 \(j=1\)。
\While{\(r_j\ne0\)}
\State 求 \(q_j,r_{j+1}\)，使 \(r_{j-1}=q_jr_j+r_{j+1}\)，且 \(r_{j+1}=0\) 或 \(\delta(r_{j+1})<\delta(r_j)\)。
\State 置 \(u_{j+1}=u_{j-1}-q_ju_j\)、\(v_{j+1}=v_{j-1}-q_jv_j\)。
\State 令 \(j\leftarrow j+1\)。
\EndWhile
\State \Return \(d=r_{j-1},\ u=u_{j-1},\ v=v_{j-1}\)。
\end{algorithmsteps}
\end{algorithm}

\subsection{整数环及一元多项式环}
\label{b1:subsec:1102:02}

在 \(\mathbb Z\) 中，取 \(\delta(n)=\lvert n\rvert\)。对 \(b\ne0\)，普通整数除法可取 \(0\leq r<\lvert b\rvert\)，故它是 Euclid 函数。以 \(252,198\) 为例，四次除法所取的商依次是 \(q_1=1,q_2=3,q_3=1,q_4=2\)：
\[
252=198+54,\qquad 198=3\cdot54+36,\qquad
54=36+18,\qquad36=2\cdot18.
\]
按\eqref{b1:eq:extended-euclid-recurrence}同步更新系数，得到下表；每行都满足 \(r_j=252u_j+198v_j\)。
\begin{center}
\begin{tabular}{rrrr}
\toprule
\(j\) & \(r_j\) & \(u_j\) & \(v_j\) \\
\midrule
\(0\) & \(252\) & \(1\) & \(0\) \\
\(1\) & \(198\) & \(0\) & \(1\) \\
\(2\) & \(54\) & \(1\) & \(-1\) \\
\(3\) & \(36\) & \(-3\) & \(4\) \\
\(4\) & \(18\) & \(4\) & \(-5\) \\
\(5\) & \(0\) & \(-11\) & \(14\) \\
\bottomrule
\end{tabular}
\end{center}
最后非零余数为 \(18\)；同一结果也可由回代读出：
\begin{equation}\label{b1:eq:integer-bezout-example}
18=54-36=4\cdot54-198=4\cdot252-5\cdot198.
\end{equation}
这里不仅算出最大公因子，也获得了以后解线性同余式所需的整数系数。

\begin{proposition}[多项式带余除法]\label{b1:prop:field-poly-division}
设 \(F\) 为域，\(f,g\in F[x]\)，且 \(g\ne0\)。存在唯一的 \(q,r\in F[x]\)，使
\[
f=gq+r,\qquad r=0\ \text{或}\ \deg r<\deg g.
\]
因此 \(F[x]\) 关于非零多项式上的次数函数为 Euclid 整环。
\end{proposition}
\begin{proof}
若 \(f=0\) 或 \(\deg f<\deg g\)，取 \(q=0,r=f\)。其余情形设 \(\deg f=n\)、\(\deg g=m\leq n\)，首项系数分别为 \(a,b\)。因为 \(F\) 是域且 \(b\ne0\)，可以减去 \(ab^{-1}x^{n-m}g\)，把 \(f\) 的最高次项消去。所得
\[
f_1=f-ab^{-1}x^{n-m}g
\]
为零或次数严格小于 \(n\)。对次数作归纳，写成 \(f_1=gq_1+r\)，其中 \(r\) 满足要求，则
\(q=ab^{-1}x^{n-m}+q_1\) 给出所求分解。次数是非负整数，故消项步骤只能进行有限次。

若又有 \(f=gq'+r'\) 且余数满足相同次数限制，则
\(g(q-q')=r'-r\)。若 \(q-q'\ne0\)，由\cref{b1:prop:polynomial-degree}，非零多项式的次数相加，使左端次数至少为 \(\deg g\)，而右端为零或次数小于 \(\deg g\)，矛盾。因此 \(q=q'\)，再得 \(r=r'\)。
\end{proof}

域假设在首项系数 \(b\) 的求逆处使用。例如在 \(\mathbb Z[x]\) 中，\(x\) 无法除以常数 \(2\) 得到次数小于 \(0\) 的余数：这种余数只能为零，而 \(x\notin2\mathbb Z[x]\)。这不否认某些特殊除数可做除法，例如首一多项式的首项系数就是单位。

\ProseParagraphBreak
在 \(\mathbb Q[x]\) 中取 \(f=x^3-2x+1\)、\(g=x^2-1\)。逐项相减得到
\[
f=xg+(-x+1),\qquad g=(-x-1)(-x+1).
\]
所以首一最大公因子为 \(x-1\)，且 \(x-1=xg-f\)。从最后非零余数 \(-x+1\) 改为首一者，需要同时将 Bézout 系数乘以 \(-1\)。这种归一化不影响生成的主理想。

\subsection{Gauss 整数的例子}
\label{b1:subsec:1102:03}

考虑复数的子环
\[
\mathbb Z[i]=\{\,a+bi\mid a,b\in\mathbb Z\,\},\qquad i^2=-1,
\]
称为\term[gauss-zhengshu]{Gauss 整数环}[ring of Gaussian integers]。它是整环，因为它嵌入 \(\mathbb C\)。定义范数
\[
N(a+bi)\coloneqq a^2+b^2.
\]
由 \(N(z)=z\overline z\) 得 \(N(zw)=N(z)N(w)\)，非零元素的范数为正整数。一个元素是单位当且仅当范数为 \(1\)：必要性由逆元与范数的乘法性得到；充分性由 \(z\overline z=1\) 得到。因此单位恰为 \(\pm1,\pm i\)。
\symidx{\(\mathbb Z[i]\)：Gauss 整数环}

若 \(\alpha\mid\beta\)，则 \(N(\alpha)\mid N(\beta)\)；反向推断一般不成立。即使范数相同，两个元素也未必相伴。例如 \(N(2+i)=N(2-i)=5\)，但
\[
\frac{2-i}{2+i}=\frac{3-4i}{5}\notin\mathbb Z[i],
\]
所以 \(2+i\nmid2-i\)，二者也不相伴。范数为整数素数足以证明一个 Gauss 整数不可约，却不是必要条件。例如 \(3\) 在 \(\mathbb Z[i]\) 中不可约（见\cref{b1:ex:11-gaussian-rational-prime}），但 \(N(3)=9\)。

\begin{proposition}\label{b1:prop:gaussian-euclidean}
设
\[
R=\mathbb Z[i]=\{\,a+bi\mid a,b\in\mathbb Z\,\},\qquad
N(a+bi)=a^2+b^2.
\]
则 Gauss 整数环 \(R\) 关于 \(R\setminus\{0\}\) 上的范数 \(N\) 是 Euclid 整环。
\end{proposition}
\begin{proof}
对 \(\alpha,\beta\in\mathbb Z[i]\)、\(\beta\ne0\)，在复数域中写 \(\alpha/\beta=s+ti\)。分别选最近的整数 \(m,n\)，使
\(\lvert s-m\rvert\leq1/2\)、\(\lvert t-n\rvert\leq1/2\)。置 \(q=m+ni\)、\(r=\alpha-\beta q\)，则
\[
N(r)=N(\beta)\bigl((s-m)^2+(t-n)^2\bigr)
\leq\tfrac12N(\beta)<N(\beta).
\]
因 \(q,r\in\mathbb Z[i]\)，这正是所需带余除法；若 \(r=0\)，也满足定义。
\end{proof}

例如 \(5=(2-i)(2+i)\)，所以 \(2+i\) 整除 \(5\)。又 \(N(2+i)=5\) 为整数素数，任何分解 \(2+i=\alpha\beta\) 都使两个非负整数范数的积为 \(5\)，其中一个必为 \(1\)；故 \(2+i\) 不可约。在 \(\mathbb Z\) 中不可约的 \(5\)，进入 \(\mathbb Z[i]\) 后已经分裂。这提示我们：讨论不可约性之前，工作环必须说清楚。

% !TeX root = ../../Book1.tex
% 纲要标题：### 11.3 主理想整环
% 纲要位置：计划纲要.md，第 428 行。

\section{主理想整环}
\label{b1:sec:1103}

% 纲要内容（仅作写作计划，尚非教材正文）：
%
% * 理想主性与素元性质
% * 主理想整环是唯一分解整环
% * 非 Euclid 的主理想整环作为选读例子
%

Euclid 算法提供了求最大公因子的方法，但它还有一个不依赖具体计算的后果：每个理想都能用一个元素生成。把这一结论单独提出，许多分解论证明便不再需要一个指定的除法算法。最后的例子将说明，这样做确实扩大了研究范围。

\subsection{理想主性与素元性质}
\label{b1:subsec:1103:01}

\begin{definition}\label{b1:def:pid}
设 \(R\) 是交换整环。若对每个理想 \(I\subseteq R\)，都存在 \(d\in R\) 使 \(I=(d)=\{\,dr\mid r\in R\,\}\)，则称 \(R\) 为\term[zhulixiangzhenghuan]{主理想整环}[principal ideal domain]，简称 PID。
\end{definition}

定义包含零理想 \((0)\) 和单位理想 \((1)\)。域只有这两个理想，因而是主理想整环；“整环”这一要求不能删去后仍沿用同一个概念。

\ProseParagraphBreak
要从 Euclid 除法得到理想主性，可以把除法用在理想内部：选定一个非零元素 \(d\) 后，任何不被 \(d\) 整除的理想元素都会产生一个非零余数，而这个余数仍在理想中。因此，若 \(d\) 的 Euclid 值已经最小，就不能再有这样的元素，\(d\) 必须整除整个理想。

\begin{theorem}\label{b1:thm:euclidean-pid}
设 \(R\) 为 Euclid 整环，则 \(R\) 是主理想整环。
\end{theorem}
\begin{proof}
设 \(R\) 带 Euclid 函数 \(\delta\)，\(I\) 是非零理想。从 \(I\setminus\{0\}\) 中取 \(d\)，使 \(\delta(d)\) 最小。对任意 \(a\in I\)，写 \(a=qd+r\)，其中 \(r=0\) 或 \(\delta(r)<\delta(d)\)。但 \(r=a-qd\in I\)，最小性排除了非零余数，所以 \(a\in(d)\)。反向包含由 \(d\in I\) 得出，故 \(I=(d)\)。零理想本来就由 \(0\) 生成。
\end{proof}

\begin{proposition}\label{b1:prop:pid-bezout}
设 \(R\) 为主理想整环，\(a,b,d\in R\)。则 \(a,b\) 有最大公因子，且
\[
(a,b)=(d)\quad\Longleftrightarrow\quad d\text{ 是 }a,b\text{ 的最大公因子}.
\]
此时存在 \(u,v\in R\) 使 \(d=ua+vb\)。
\end{proposition}
\begin{proof}
主性给出 \((a,b)=(d_0)\)。其中 \(a,b\in(d_0)\) 表明 \(d_0\) 是公因子，而 \(d_0\in(a,b)\) 给出线性组合表达式；\cref{b1:prop:bezout-gcd}于是说明它是最大公因子。任一其他最大公因子 \(d\) 都与 \(d_0\) 相伴，故\cref{b1:prop:associate-ideals}给出 \((d)=(d_0)=(a,b)\)，再由 \(d\in(a,b)\) 得到相应 Bézout 系数。这也包含 \(a=b=0\) 的情形。
\end{proof}

\begin{proposition}\label{b1:prop:pid-irreducible-prime}
在主理想整环 \(R\) 中，每个不可约元 \(p\) 都是素元，且 \((p)\) 是极大理想。因此每个非零素理想都是极大理想。
\end{proposition}
\begin{proof}
设 \((p)\subseteq I\subseteq R\)。写 \(I=(d)\)，则 \(p=dc\)。不可约性使 \(d\) 为单位或 \(c\) 为单位；前者给 \(I=R\)，后者给 \(I=(p)\)。又 \(p\) 非单位，所以 \((p)\) 是极大理想。由\cref{b1:prop:prime-max-quotient}，\(R/(p)\) 是域，因而是整环，\((p)\) 为素理想；再用\cref{b1:prop:prime-implies-irreducible}的素理想判别得 \(p\) 为素元。

若 \(P\ne(0)\) 是素理想，写 \(P=(p)\)。因 \(P\) 非零且真，\(p\) 非零非单位；同一判别说明 \(p\) 为素元，再由\cref{b1:prop:prime-implies-irreducible}得它不可约。刚才的论证便说明 \(P\) 极大。
\end{proof}

\subsecref{b1:subsec:1101:03}已把不可约性解释为“在真主理想中极大”。主理想整环中每个理想都为主理想，所以这一极大性涵盖了全部真理想。上述证明中的推理可写为
\[
p\text{ 不可约}
\quad\Longrightarrow\quad
(p)\text{ 为极大理想}
\quad\Longrightarrow\quad
(p)\text{ 为素理想}
\quad\Longrightarrow\quad
p\text{ 为素元}.
\]
“每个非零素理想都是极大理想”中的非零条件也有作用。在整环中，\((0)\) 总是素理想；但在主理想整环 \(\mathbb Z\) 中，有 \((0)\subsetneq(2)\subsetneq\mathbb Z\)，所以 \((0)\) 不是极大理想。

\ProseParagraphBreak
例如对域 \(F\)，\cref{b1:prop:field-poly-division}与\cref{b1:thm:euclidean-pid}给出 \(F[x]\) 是主理想整环。若 \(p(x)\) 不可约，商 \(F[x]/(p)\) 因而是域。这个结论会在构造方程根所在的域时发挥作用；它所依赖的是理想极大性，不是预先知道根的存在。

\subsection{主理想整环的唯一分解性}
\label{b1:subsec:1103:02}

理想主性还将保证元素的不可约分解存在且唯一。这里先定义唯一分解整环，以陈述主理想整环的这一结构性质；\secref{b1:sec:1104}将脱离主理想性，单独研究唯一分解本身及其判别。

\begin{definition}\label{b1:def:ufd}
设 \(R\) 是交换整环。若每个非零非单位都能写成有限个不可约元的乘积，并且同一非零非单位 \(a\in R\) 的任意两种表达式
\[
a=u p_1\cdots p_m=v q_1\cdots q_n
\]
中，\(u,v\) 为单位、\(p_i,q_j\) 为不可约元时，必有 \(m=n\)，且可重新排列 \(q_j\)，使每个 \(p_i\) 与 \(q_i\) 相伴，则称 \(R\) 为\term[weiyifenjiezhenghuan]{唯一分解整环}[unique factorization domain]，简称 UFD。
\end{definition}

单位允许看成空乘积再乘一个单位；零不属于分解定理的对象。唯一性既不固定因子的次序，也不固定各因子相伴类中的代表。因此 \(6=2\cdot3=(-2)(-3)\) 并不是两种不同的素因数分解。

\ProseParagraphBreak
一般理想的并未必是理想，甚至未必对加法封闭。例如 \(2\mathbb Z\cup3\mathbb Z\) 包含 \(2,3\)，却不包含 \(2+3=5\)。理想升链的并有所不同：属于不同项的两个元素总能一起放入其中较大的一项，这使理想的封闭条件得以保留。

\begin{lemma}\label{b1:lem:pid-ascending-chain}
设 \(R\) 为主理想整环。则 \(R\) 的每条理想升链
\(I_1\subseteq I_2\subseteq\cdots\) 都从某项开始恒定。
\end{lemma}
\begin{proof}
令 \(I=\bigcup_{j\geq1}I_j\)。若 \(x\in I_r,y\in I_s\)，升链条件使两者都属于 \(I_{\max(r,s)}\)，故 \(x-y\) 仍在 \(I\)；任一元素乘环中元素也留在原来所属的 \(I_j\)。因此 \(I\) 是理想，写成 \((d)\)。生成元 \(d\) 属于某个 \(I_m\)，所以 \(I=(d)\subseteq I_m\subseteq I\)，从而 \(I_m=I\)，所有 \(j\geq m\) 也满足 \(I_j=I\)。
\end{proof}

分解元素与扩大主理想正好方向相反。若非零元素满足 \(a_1=a_2d_1\)，且 \(d_1\) 非单位，则 \((a_1)\subsetneq(a_2)\)。因此，若一个元素总能继续作非平凡分解，就会产生严格上升的主理想链。下面用升链终止来排除无限继续分解的情形，再用不可约元的素性比较两种分解。

\begin{theorem}\label{b1:thm:pid-ufd}
设 \(R\) 为主理想整环，则 \(R\) 是唯一分解整环。
\end{theorem}
\begin{proof}
先证分解存在。若存在不能分解成有限个不可约元乘积的非零非单位 \(a_1\)，则 \(a_1\) 本身不可为不可约元，故 \(a_1=b_1c_1\)，其中两个因子都非单位。至少有一个因子仍不能分解，否则合并两边分解便得到 \(a_1\) 的分解。把这个因子记为 \(a_2\)，另一个记为 \(d_1\)，于是 \(a_1=a_2d_1\)。有 \((a_1)\subsetneq(a_2)\)：若相等，\cref{b1:prop:associate-ideals}使 \(a_1,a_2\) 相伴，再由非零消去律迫使 \(d_1\) 为单位，矛盾。重复选择会得到严格升链
\((a_1)\subsetneq(a_2)\subsetneq\cdots\)，与\cref{b1:lem:pid-ascending-chain}矛盾。因此分解存在。

再证唯一性。设 \(u p_1\cdots p_m=v q_1\cdots q_n\)。由\cref{b1:prop:pid-irreducible-prime}，不可约元 \(p_1\) 是素元；它不整除单位 \(v\)，对乘积逐次应用素性，必整除某个 \(q_j\)。不可约性使 \(q_j=p_1w\)，其中 \(w\) 为单位。把这个因子移至首位，消去非零的 \(p_1\)，得到两个较短分解的等式，单位因子相应并入 \(v\)。如此归纳，配对因子都相伴。若一边先用尽而另一边仍有不可约因子，就会使若干非单位之积成为单位；但乘积有逆元时，每个因子也有逆元，矛盾。因此两边同时用尽，\(m=n\)，唯一性得证。
\end{proof}

这一证明的两部分使用不同的条件：
\[
\begin{array}{c@{\qquad}c}
\text{分解存在性}&\text{分解唯一性}\\
\text{理想升链终止}&\text{不可约元都是素元}
\end{array}
\]
由此前的具体 Euclid 函数，\(\mathbb Z\)、\(F[x]\) 和 \(\mathbb Z[i]\) 均为唯一分解整环。不过，主性证明中的最小值选择只断言某个生成元存在；若没有可执行的除法步骤，它本身还不是计算该生成元的算法。

\OptionalSubsection{非 Euclid 的主理想整环}{b1:subsec:1103:03}

要说明“Euclid 整环 \(\Rightarrow\) 主理想整环”的逆命题不成立，需要对同一个环完成两件事：先证明每个理想都是主理想，再排除所有可能的 Euclid 函数。前一部分将利用格点估计构造允许倍乘的较弱余数；后一部分则从任意 Euclid 函数推出一个不可能的有限商环。设
\[
\omega=\frac{1+\sqrt{-19}}2,\qquad
R=\mathbb Z[\omega]\subseteq\mathbb C.
\]
因为 \(\omega^2=\omega-5\)，每个元素唯一写成 \(a+b\omega\)，\(a,b\in\mathbb Z\)。定义
\[
N(a+b\omega)=\lvert a+b\omega\rvert^2
=a^2+ab+5b^2
=\Bigl(a+\frac b2\Bigr)^2+\frac{19b^2}4.
\]
范数是非负整数且乘法相容，非零元素的范数是正整数。若 \(u\) 是单位，则 \(N(u)N(u^{-1})=N(1)=1\)，所以 \(N(u)=1\)。由最后一个表达式，范数等于 \(1\) 时只能有 \(b=0,a=\pm1\)，而 \(\pm1\) 确实可逆，因此 \(R\) 的单位恰为 \(\pm1\)。

\ProseParagraphBreak
比较下面两种余数要求，其中 \(\beta\ne0\)。左边是以范数为度量的普通带余除法，右边只针对 \(\alpha\notin(\beta)\)：
\[
\begin{array}{c@{\qquad}c}
\text{普通带余除法}&\text{允许倍乘的较弱条件}\\
\alpha=d\beta+r&c\alpha=d\beta+r\\
r=0\ \text{或}\ N(r)<N(\beta)&0<N(r)<N(\beta)
\end{array}
\]
右边的乘数 \(c\in R\) 可以依赖于 \(\alpha/\beta\)，所以所得 \(r\) 是 \(c\alpha\) 的余数，而不是 \(\alpha\) 的余数。这一条件不提供普通带余除法，却足以证明理想主性：若非零理想 \(I\) 中的 \(\beta\) 具有最小正范数，而 \(\alpha\in I\setminus(\beta)\)，则 \(r=c\alpha-d\beta\) 仍在 \(I\) 中，其正范数严格更小，产生矛盾。令 \(z=\alpha/\beta\notin R\)，右边的不等式就化为 \(0<\lvert cz-d\rvert^2<1\)。下面的格点引理构造这样的 \(c,d\)，并保证余数非零。

\begin{lemma}\label{b1:lem:nineteen-small-combination}
令
\[
\omega=\frac{1+\sqrt{-19}}2,\qquad R=\mathbb Z[\omega]\subseteq\mathbb C.
\]
对任意 \(z\in\mathbb C\setminus R\)，存在 \(c,d\in R\)，使
\(0<\lvert cz-d\rvert^2<1\)。
\end{lemma}
\begin{proof}
由 \(R=\mathbb Z+\mathbb Z\omega\) 及 \(\operatorname{Im}\omega=\sqrt{19}/2\)，减去适当整数倍的 \(\omega\) 可以把虚部移入一个长度为 \(\sqrt{19}/2\)、以零为中心的区间；再减去整数，把实部移入长度为 \(1\)、以零为中心的区间。因此可取 \(\rho\in R\)，使调整后的 \(z_0=z-\rho=x+yi\) 落入基本矩形
\[
\lvert x\rvert\leq\frac12,\qquad
\lvert y\rvert\leq\frac{\sqrt{19}}4.
\]
由于 \(z\notin R\)，调整后的 \(z_0\) 仍不在 \(R\) 中，否则 \(z=z_0+\rho\in R\)。特别地 \(z_0\ne0\)。若 \(\lvert z_0\rvert<1\)，取 \(c=1,d=\rho\)，就有 \(0<\lvert cz-d\rvert^2=\lvert z_0\rvert^2<1\)，已经得到结论。

\cref{b1:fig:nineteen-lattice}中矩形与单位圆之间的上下两片区域，正是基本矩形内仍须处理 \(\lvert z_0\rvert\geq1\) 的点。图中取 \(z_0=\omega/2\)：从 \(z_0\) 到 \(\omega\) 是倍乘，从 \(\omega\) 回到原点才是减去格点 \(\omega\)。

\begin{figure}[htbp]
\centering
\begin{tikzpicture}[x=1.55cm,y=1.55cm,>=stealth]
  \def\h{1.08972474} % sqrt(19)/4
  \def\s{0.86602540} % sqrt(3)/2
  \fill[color=AlgebraBlueLight] (-.5,\s)--(-.5,\h)--(.5,\h)--(.5,\s)
    arc[start angle=60,end angle=120,radius=1]--cycle;
  \fill[color=AlgebraBlueLight] (-.5,-\s)--(-.5,-\h)--(.5,-\h)--(.5,-\s)
    arc[start angle=-60,end angle=-120,radius=1]--cycle;
  \draw[->,color=AlgebraGray] (-1.2,0)--(1.25,0) node[right] {\(\operatorname{Re}z\)};
  \draw[->,color=AlgebraGray] (0,-1.26)--(0,2.51) node[above] {\(\operatorname{Im}z\)};
  \draw[dashed,color=AlgebraGray] (0,0) circle (1);
  \draw[thick,color=AlgebraBlue] (-.5,-\h) rectangle (.5,\h);
  \fill (0,0) circle (1.2pt) node[below left] {\(0\)};
  \fill[color=AlgebraBlue] (.25,\h) circle (1.6pt);
  \node[anchor=east,color=AlgebraBlue] at (.17,1.30) {\(z_0=\omega/2\)};
  \fill[color=AlgebraBlue] (.5,2*\h) circle (1.6pt);
  \node[anchor=west,color=AlgebraBlue] at (.55,2*\h) {\(2z_0=\omega\)};
  \draw[->,thick,color=AlgebraBlue] (.28,\h+.08)--(.48,2*\h-.09)
    node[pos=.53,left] {\(\times2\)};
  \draw[->,thick,color=AlgebraGray] (.55,2*\h-.03)
    .. controls (1.68,2.02) and (1.72,.38) .. (.06,.04);
  \node[color=AlgebraGray] at (1.38,1.16) {\(-\omega\)};
\end{tikzpicture}
\caption[基本矩形与单位圆]{格点引理中的基本矩形与单位圆（虚线）。阴影为矩形内满足 \(\lvert z_0\rvert\geq1\) 的上下两片区域；箭头分别表示倍乘与减去格点，图中展示的是零余数实例。}
\label{b1:fig:nineteen-lattice}
\end{figure}

以下处理基本矩形内仍有 \(\lvert z_0\rvert\geq1\) 的情形。由 \(x^2\leq1/4\)，有 \(\lvert y\rvert\geq\sqrt3/2\)。对 \(2z_0\) 再作格点调整：减去虚部与 \(y\) 同号的 \(\omega\) 或 \(-\omega\)，再减去整数使实部绝对值不超过 \(1/2\)，得到 \(w=2z_0-d_0\)，且
\[
\lvert\operatorname{Im}w\rvert
\leq\frac{\sqrt{19}}2-\sqrt3<\frac12,
\qquad \lvert w\rvert^2<\frac12<1.
\]
其中严格不等式可由 \(\sqrt{19}<22/5\)、\(\sqrt3>17/10\) 核对。若 \(w\ne0\)，取 \(c=2,d=d_0+2\rho\)，便有 \(cz-d=2z_0-d_0=w\)，即得结论。

最后处理 \(w=0\)。此时倍乘 \(c=2\) 恰好把 \(z_0\) 送到格点，所得余数是零，不能用它反驳最小正范数。我们改换乘数 \(c\)，目标是制造余数 \(1/2\)，从而同时保证余数非零与模平方小于 \(1\)。令 \(\alpha=2z_0\in R\)；由于 \(z_0\notin R\)，\(\alpha\notin2R\)。

按 \(a+b\omega\) 的两个整数系数的奇偶性，商环 \(R/2R\) 恰有四个剩余类，分别由 \(0,1,\omega,1+\omega\) 代表。三个非零类的逆元分别由 \(1,1+\omega,\omega\) 代表，因为
\(\omega(1+\omega)=2\omega-5\equiv1\pmod{2R}\)。这说明 \(R/2R\) 是四元素域；当前所需的正是非零类可逆。\(\alpha\) 的类非零，故可取 \(c\in R\) 及 \(d_1\in R\)，使 \(c\alpha=1+2d_1\)。于是 \(cz_0-d_1=1/2\)，其模平方为 \(1/4\)。取 \(d=d_1+c\rho\)，便有 \(cz-d=cz_0-d_1=1/2\)，就得到原命题。

例如 \(z_0=\omega/2\) 位于图中的上方阴影区域，因为 \(\lvert z_0\rvert^2=5/4>1\)，但 \(2z_0-\omega=0\)。这时可直接取 \(c=1+\omega\)、\(d_1=\omega-3\)；由 \(\omega^2=\omega-5\) 得
\[
(1+\omega)\frac{\omega}{2}-(\omega-3)=\frac12.
\]
所以零余数分支也确实产生非零且模平方小于 \(1\) 的线性组合。
\end{proof}

\begin{proposition}\label{b1:prop:noneuclidean-pid}
令 \(\omega=(1+\sqrt{-19})/2\)，并令 \(R=\mathbb Z[\omega]\subseteq\mathbb C\)。则 \(R\) 是主理想整环，但不存在任何使 \(R\) 成为 Euclid 整环的 Euclid 函数。
\end{proposition}
\begin{proof}
设 \(I\ne(0)\) 为理想，取 \(\beta\in I\setminus\{0\}\)，使正整数 \(N(\beta)\) 最小。若某个 \(\alpha\in I\) 不属于 \((\beta)\)，则 \(z=\alpha/\beta\notin R\)。由\cref{b1:lem:nineteen-small-combination}取 \(c,d\in R\)，使 \(0<\lvert cz-d\rvert^2<1\)。于是非零元素 \(c\alpha-d\beta\in I\) 满足
\[
0<N(c\alpha-d\beta)
=N(\beta)\lvert cz-d\rvert^2<N(\beta),
\]
与最小性矛盾。所以 \(I=(\beta)\)，零理想也为主理想，得证主性。

现在排除任意 Euclid 函数。由于 \(N(2)=4\ne1\)，\(2\) 是非零非单位，因此所有非零非单位组成的集合非空。假设存在 Euclid 函数 \(\delta\)，可在这个集合中取 \(b\)，使 \(\delta(b)\) 最小。对任意 \(a\in R\) 作除法 \(a=qb+r\)；若 \(r\ne0\)，则 \(\delta(r)<\delta(b)\)。最小性使 \(r\) 不可能是非单位，因而只能是单位，即 \(1\) 或 \(-1\)。

因此每个模 \((b)\) 的剩余类都由 \(0,1,-1\) 之一代表，\(R/(b)\) 至多有三个元素。另一方面，\(b\) 非单位意味着 \(1\notin(b)\)，所以商环中的 \(1\ne0\)，至少有两个元素。合起来得到
\[
2\leq\lvert R/(b)\rvert\leq3.
\]
故它的大小只能为 \(2\) 或 \(3\)。

令 \(S=R/(b)\)，并令 \(p=\lvert S\rvert\in\{2,3\}\)。由\cref{b1:thm:lagrange}，\(1_S\ne0_S\) 在加法群中的阶是 \(p\)，所以 \(1_S\) 生成整个加法群。映射
\[
\mathbb Z/p\mathbb Z\longrightarrow S,\qquad
[n]\longmapsto n1_S
\]
于是良定义且为双射；它保持加法、乘法和单位，其中乘法相容性来自 \((m1_S)(n1_S)=(mn)1_S\)。因此它是含幺环同构，即 \(S\cong\mathbb F_p\)。

然而 \(\omega\) 的像必须满足 \(t^2-t+5=0\)。模 \(2\) 时，\(t=0,1\) 的值都是 \(1\)；模 \(3\) 时，\(t=0,1,2\) 的值依次为 \(2,2,1\)。两种商中都不可能有这样的元素，矛盾。故不存在 \(\delta\)。
\end{proof}

范数本身不能作 Euclid 函数，还不足以推出环不是 Euclid 整环；必须排除所有可能的 Euclid 函数。上述证明的后半部分正是通过有限商环的必要条件完成这一点。主性证明使用的乘数 \(c\) 则足以产生理想中的更小元素，却不提供对原元素的普通带余除法。

\ProseParagraphBreak
关于这个例子，可参见 Keith Conrad 的讲义\cite[Theorems 5.18, 5.22]{ConradRemarksEuclideanDomains}：定理 5.18 提炼了非域 Euclid 整环中某个非单位模下每个剩余类可由零或单位代表的必要条件；定理 5.22 讨论本小节的二次环，其中主理想性只给出证明提纲，可与此处的格点引理作方法对照。

% !TeX root = ../../Book1.tex
% 纲要标题：### 11.4 唯一分解整环
% 纲要位置：计划纲要.md，第 434 行。

\section{唯一分解整环}
\label{b1:sec:1104}

% 纲要内容（仅作写作计划，尚非教材正文）：
%
% * 分解的存在和唯一性
% * 唯一分解不蕴含主性
% * 非唯一分解例子与理想分解的动机
%
% ---
%

主理想性给出了唯一分解的充分条件，却不是必要条件。本节把分解论自身的要求提炼出来，并考察两个相反方向的例子：\(\mathbb Z[x]\) 的元素仍能唯一分解，但有非主理想；\(\mathbb Z[\sqrt{-5}]\) 连元素的分解唯一性也会失败。前者的唯一分解性将在下一章通过多项式方法证明，后者在这里直接计算。

\subsection{分解的存在性与唯一性}
\label{b1:subsec:1104:01}

\begin{proposition}\label{b1:prop:ufd-prime}
设 \(R\) 为整环。如果 \(R\) 是唯一分解整环，则 \(R\) 中的不可约元都是素元。更一般地，\(R\) 是唯一分解整环，当且仅当它的每个非零非单位都能分解成有限个不可约元的乘积，并且每个不可约元都是素元。
\end{proposition}
\begin{proof}
先设 \(R\) 为唯一分解整环，\(p\) 不可约且 \(p\mid ab\)。若 \(a=0\) 或 \(b=0\)，结论成立；否则写 \(ab=pc\)，其中 \(c\ne0\)。把 \(a,b,c\) 分别分解为单位与不可约元的乘积，允许单位的分解含零个不可约因子。等式右侧有因子 \(p\)，唯一性要求它与左侧某个不可约因子相伴。那个因子来自 \(a\) 或 \(b\)，所以 \(p\mid a\) 或 \(p\mid b\)，即 \(p\) 为素元。

反过来，假定分解存在且不可约元都为素元。\cref{b1:thm:pid-ufd}证明唯一性的后一段，在已有两种分解后只用了不可约元的素性、整环的非零消去律及单位的基本性质，并未再用理想主性。因此同一因子配对论证适用于此，得到因子个数相同且对应因子相伴，正是\cref{b1:def:ufd}的唯一性。
\end{proof}

这个判别不能省略分解存在的条件。证明唯一性时以“取一个不可约分解”起步，只在确知每个非零非单位都能如此分解时才有意义。\cref{b1:thm:pid-ufd}为主理想整环补足存在性时，用到的是理想升链最终稳定，两个论证的职责应当分清。

设 \(a,b\) 为整环 \(R\) 的非零元素。如果 \(m\in R\) 同时被 \(a,b\) 整除，并且 \(m\) 整除 \(a,b\) 的每个公倍数，就称 \(m\) 为 \(a,b\) 的\term[zuixiaogongbeishu]{最小公倍数}[least common multiple]，记为 \(\operatorname{lcm}(a,b)\)。最小公倍数只在相伴意义下确定。
\symidx{\(\operatorname{lcm}(a,b)\)：最小公倍数}

\begin{proposition}[唯一分解整环中的最大公因子与最小公倍数]\label{b1:prop:ufd-gcd-lcm}
设 \(R\) 为唯一分解整环，\(a,b\in R\) 非零。则 \(a,b\) 有最大公因子与最小公倍数。把它们分解中涉及的不同不可约相伴类各选一个代表 \(p_1,\ldots,p_t\)，其中 \(t\geq0\) 为整数，并写成
\[
a=u\prod_{j=1}^{t}p_j^{\alpha_j},\qquad
b=v\prod_{j=1}^{t}p_j^{\beta_j},\qquad \alpha_j,\beta_j\in\mathbb Z_{\geq0},
\]
其中 \(u,v\in R^\times\)。此时最大公因子 \(d\) 和最小公倍数 \(m\) 可分别取为
\[
d=\prod_{j=1}^{t}p_j^{\min(\alpha_j,\beta_j)},\qquad
m=\prod_{j=1}^{t}p_j^{\max(\alpha_j,\beta_j)}.
\]
\end{proposition}
\begin{proof}
唯一分解说明，对于非零元素 \(r,s\)，条件 \(r\mid s\) 等价于 \(r\) 中每个不可约因子的指数不超过它在 \(s\) 中的指数：必要性由 \(s=rc\) 及分解唯一性得到，充分性由指数相减构造商得到。因此公因子中各指数的上界是两个指数的较小者，公倍数中各指数的下界是两个指数的较大者，分别给出所列 \(d,m\)。公倍数为 \(0\) 时，\(m\mid0\) 也成立。
\end{proof}

对上一命题选出的 \(d,m\)，乘积 \(dm\) 与 \(ab\) 相伴；适当调整单位后可令二者相等。零元素的情况另定为 \(\gcd(a,0)\) 与 \(a\) 相伴、\(\operatorname{lcm}(a,0)=0\)。

\ProseParagraphBreak
例如在 \(\mathbb Z[i]\) 中，\(2=-i(1+i)^2\)，所以 \(2\) 已不再是不可约元。若 \(a=(1+i)^3(2+i)\)、\(b=(1+i)(2+i)^2\)，则 \(1+i\) 与 \(2+i\) 的范数分别为整数素数 \(2,5\)，两者不可约且不相伴。因而可取 \(\gcd(a,b)=(1+i)(2+i)\)，最小公倍数则为 \((1+i)^3(2+i)^2\)。

\subsection{唯一分解不蕴含主性}
\label{b1:subsec:1104:02}

\cref{b1:prop:integer-poly-nonprincipal}已经说明：\(\mathbb Z[x]\) 的理想 \((2,x)\) 不为主理想。这里的障碍是 \(2,x\) 虽然没有共同的不可约因子，却不能用多项式系数的线性组合得到 \(1\)。

\ProseParagraphBreak
另一方面，\(\mathbb Z[x]\) 确实是唯一分解整环。这个断言是下一章\cref{b1:thm:polynomial-ufd}的一个应用：那里将证明只要 \(R\) 是唯一分解整环，\(R[x]\) 也是；取本章已证的 \(R=\mathbb Z\) 即得。此处把它作为前瞻例子，不用它证明本章其他结论。与已经完成证明的 \(\mathbb Z[\omega]\) 一起，这两个例子说明三个条件之间是严格蕴含：
\[
\text{Euclid 整环}\ \Longrightarrow\
\text{主理想整环}\ \Longrightarrow\
\text{唯一分解整环}.
\]
左边逆向的反例由\cref{b1:prop:noneuclidean-pid}给出，右边逆向的反例则在下一章完成最后的唯一分解性验证。

\ProseParagraphBreak
在 \(\mathbb Z[x]\) 中，逐项模 \(2\) 的商映射给出 \(\mathbb Z[x]/(2)\cong\mathbb F_2[x]\)。右边是整环而不是域，所以 \(2\) 是素元，\((2)\) 却不是极大理想；事实上 \((2)\subsetneq(2,x)\subsetneq\mathbb Z[x]\)。下一章\cref{b1:thm:polynomial-ufd}还将证明 \(\mathbb Z[x]\) 是唯一分解整环；因而这一例子也表明，不可约元都是素元并不能保证每个非零素理想都极大。

\ProseParagraphBreak
在唯一分解整环中，对非零 \(a,b\)，公倍数恰为最小公倍数的倍数，因此 \((a)\cap(b)=(m)\) 仍是主理想；而和 \((a)+(b)\) 未必为主理想。最大公因子刻画公因子之间的整除关系，却未必能表示为原来两个元素的线性组合。将“有最大公因子”与“能写 Bézout 关系”分开，能避免把主理想整环的结论错误推广到全部唯一分解整环。

\subsection{非唯一分解与理想分解}
\label{b1:subsec:1104:03}

令 \(s=\sqrt{-5}\)，考虑整环 \(R=\mathbb Z[s]\subseteq\mathbb C\)。它的元素唯一写成 \(a+bs\)，范数
\[
N(a+bs)=a^2+5b^2
\]
是乘法相容的非负整数，单位恰为 \(\pm1\)。不存在范数为 \(2\) 或 \(3\) 的元素：若 \(b\ne0\)，范数至少为 \(5\)；若 \(b=0\)，整数平方又不可能等于 \(2,3\)。

\begin{proposition}\label{b1:prop:sqrt-minus-five-factorization}
令 \(s=\sqrt{-5}\)，并令 \(R=\mathbb Z[s]\subseteq\mathbb C\)。则等式
\[
6=2\cdot3=(1+s)(1-s)
\]
给出两种不等价的不可约分解。因此 \(R\) 不是唯一分解整环，且不可约元 \(2\) 不是素元。
\end{proposition}
\begin{proof}
元素 \(2,3,1+s,1-s\) 的范数依次为 \(4,9,6,6\)。若 \(2\) 分解成两个非单位，两因子的范数只能为 \(2,2\)，但范数 \(2\) 不存在，故 \(2\) 不可约。对 \(3\)，唯一的非平凡范数分解为 \(3\cdot3\)，同样不可能。对 \(1\pm s\)，非平凡范数分解只能为 \(2\cdot3\)，也不可能。因此四个因子都不可约。

单位的范数为 \(1\)，相伴元素范数相同，所以范数为 \(4\) 的 \(2\) 不可能与范数为 \(6\) 的任一 \(1\pm s\) 相伴。两种分解因而不等价。此外，\(2\mid(1+s)(1-s)=6\)，却不整除任一 \(1\pm s\)，因为它们的整数常数项为奇数。故 \(2\) 不满足素性。
\end{proof}

元素分解失效以后，可以尝试让理想承担因子的作用。本例的
\[
I=(2,1+s)
\]
不是主理想，却满足 \(I^2=(2)\)。先看非主性。映射
\(R\rightarrow\mathbb F_2\)、\(a+bs\mapsto\overline{a+b}\)
保持乘法，因为 \(s^2=-5\) 在模 \(2\) 下对应 \(1^2=1\)；它把 \(I\) 送到零而把 \(1\) 送到 \(1\)，故 \(I\ne R\)。若 \(I=(d)\)，则 \(d\mid2\)。由 \(2\) 不可约，\(d\) 或为单位，或与 \(2\) 相伴；前者与 \(I\ne R\) 矛盾，后者又使 \(1+s\in(2)\)，也矛盾。

\ProseParagraphBreak
再算理想平方。\(I^2\) 由
\[
4,\quad 2+2s,\quad (1+s)^2=-4+2s
\]
生成，三者都属于 \((2)\)。反向包含由
\[
2=(2+2s)-(-4+2s)-4\in I^2
\]
得到，所以 \(I^2=(2)\)。元素 \(2\) 不可约，但它生成的理想仍有这样的真分解。\cref{b1:ex:11-ideal-splitting-three}将继续核对 \(J=(3,1+s)\)、\(K=(3,1-s)\) 满足 \(JK=(3)\)、\(IJ=(1+s)\)、\(IK=(1-s)\)。于是原来的两种元素分解在理想层面汇合为
\[
(6)=(2)\cdot(3)=I^2JK=(IJ)(IK)
=\bigl((1+s)R\bigr)\cdot\bigl((1-s)R\bigr).
\]
最后一项明确写出两个主理想的乘积，以区别于前面的元素等式 \(6=(1+s)(1-s)\)。这些等式只核对当前例子；要得到一般的理想分解定理，还须对环施加适当条件。这是后续数论与交换代数需要研究的问题。


\begin{chaptersummary}
素元必不可约，反向蕴含需要额外条件。Euclid 除法使连续非零余数的 Euclid 函数值严格下降，故迭代过程必定终止，最后一个非零余数给出最大公因子。若已给出可执行的求商余程序，还可同步更新系数，计算 Bézout 关系。对非零理想选取 Euclid 函数值最小的非零元素，可证明它生成该理想。在主理想整环中，理想升链最终稳定保证不可约分解存在，不可约元的素性保证分解唯一。

\ProseParagraphBreak
唯一分解整环中的最大公因子可以通过因子指数求得，却不一定能写成 Bézout 关系。\(\mathbb Z[x]\) 的理想 \((2,x)\) 展示了这种差别。\(\mathbb Z[\sqrt{-5}]\) 中不可约但不素的 \(2\)，则说明元素的唯一分解可能彻底失效；相应的主理想仍可分解为非主理想的乘积。下一章将通过容度与本原多项式，解决多项式环何时保留唯一分解的问题。
\end{chaptersummary}

\BookExercises

\begin{exercise}\label{b1:ex:11-gaussian-gcd}
在 \(\mathbb Z[i]\) 中求 \(a=7+5i\)、\(b=3+i\) 的一个最大公因子及 Bézout 系数，并列出它们的公因子相伴类。求 \(a,b\) 的一个最小公倍数，并计算理想 \((a)+(b)\) 与 \((a)\cap(b)\)。
\end{exercise}
\begin{solution}
选择商 \(3+i\)，有
\[
a=(3+i)b+(-1-i),\qquad
b=(-2+i)(-1-i).
\]
第一步余数范数为 \(2<N(b)=10\)，第二步余数为零，故\cref{b1:prop:euclidean-algorithm}给出最大公因子 \(-1-i\)。乘以单位 \(-1\)，可取 \(d=1+i\)，相应关系为
\[
1+i=-a+(3+i)b.
\]
范数 \(N(1+i)=2\) 为整数素数，若 \(1+i\) 分解，两因子范数的积为 \(2\)，必有一个范数为 \(1\)，即为单位。所以 \(1+i\) 不可约。公因子恰为 \(d\) 的因子，只有单位类和 \(1+i\) 的相伴类，分别以 \(1,1+i\) 为代表。

由\cref{b1:prop:ufd-gcd-lcm}，可取最小公倍数
\[
m=\frac{ab}{d}=19+3i.
\]
前述 Bézout 关系给出 \((a)+(b)=(d)=(1+i)\)。理想 \((a)\cap(b)\) 中的元素恰为 \(a,b\) 的公倍数；最小公倍数的整除刻画因而给出 \((a)\cap(b)=(m)=(19+3i)\)。
\end{solution}

\begin{exercise}\label{b1:ex:11-diophantine}
求出方程 \(252u+198v=36\) 的全部整数解，并说明右端改为 \(30\) 时是否有解。
\end{exercise}
\begin{solution}
将\eqref{b1:eq:integer-bezout-example}的 Bézout 关系乘以 \(2\)，得到 \(36=8\cdot252-10\cdot198\)，故 \((u,v)=(8,-10)\) 是一个解。任意其他解满足
\[
14(u-8)+11(v+10)=0.
\]
模 \(11\) 后，\(3(u-8)\equiv0\pmod{11}\)。因 \(4\cdot3\equiv1\pmod{11}\)，得到 \(u-8=11t\)，代回得 \(v+10=-14t\)。因此全部解为
\[
u=8+11t,\qquad v=-10-14t,\qquad t\in\mathbb Z.
\]
反向代入可检验每个 \(t\) 均给出解。右端为 \(30\) 时无解，因为左端必被 \(18\) 整除，而 \(18\nmid30\)。
\end{solution}

\begin{exercise}\label{b1:ex:11-polynomial-euclid}
在 \(\mathbb F_2[x]\) 中，求
\[
f=x^4+x^3+x^2+1,\qquad g=x^3+1
\]
的首一最大公因子，并求 \(u,v\in\mathbb F_2[x]\)，使该最大公因子等于 \(uf+vg\)。须列出各次带余除法。
\end{exercise}
\begin{solution}
在特征 \(2\) 中，加法与减法相同。按\cref{b1:prop:field-poly-division}消去首项，得到
\[
\begin{aligned}
f&=(x+1)g+(x^2+x),\\
g&=(x+1)(x^2+x)+(x+1),\\
x^2+x&=x(x+1).
\end{aligned}
\]
非零余数的次数依次为 \(2,1\)，最后余数为零。由\cref{b1:prop:euclidean-algorithm}，首一最大公因子为 \(d=x+1\)。回代为
\[
\begin{aligned}
d&=g+(x+1)(x^2+x)\\
&=g+(x+1)\bigl(f+(x+1)g\bigr)\\
&=(x+1)f+x^2g.
\end{aligned}
\]
最后一步使用 \(1+(x+1)^2=x^2\)，这是当前基域特征为 \(2\) 的结果。因此可取 \(u=x+1\)、\(v=x^2\)；把这些系数误作有理数系数时，不能保留相同的运算步骤。
\end{solution}

\begin{exercise}\label{b1:ex:11-gaussian-rational-prime}
设整数素数 \(p\equiv3\pmod4\)。证明 \(p\) 在 \(\mathbb Z[i]\) 中仍为素元，并解释为什么同一论证不能用于 \(p=2\)。
\end{exercise}
\begin{solution}
若 \(p=\alpha\beta\) 且两个因子都为非单位，则两个正整数 \(N(\alpha),N(\beta)>1\) 的积为 \(p^2\)，只能同时等于 \(p\)。然而整数平方模 \(4\) 只能为 \(0,1\)，两平方之和不可能同余 \(3\)，所以不存在范数为 \(p\) 的 Gauss 整数。这证明 \(p\) 不可约。

由\cref{b1:prop:gaussian-euclidean}和\cref{b1:thm:euclidean-pid}，\(\mathbb Z[i]\) 为主理想整环；\cref{b1:prop:pid-irreducible-prime}因此把这里的不可约性提升为素性。当 \(p=2\) 时，范数 \(2\) 确实存在，例如 \(N(1+i)=2\)，并且 \(2=-i(1+i)^2\)，所以 \(2\) 在此环中并不不可约。
\end{solution}

\begin{exercise}\label{b1:ex:11-quotient-inverse}
设 \(R\) 为主理想整环，\(a\) 为非零非单位。证明 \(b+(a)\) 在 \(R/(a)\) 中可逆，当且仅当 \(a,b\) 互素。应用于 \(R=\mathbb Z[i]\)、\(a=3\)，证明该商为九元素域，并求 \(1+i+(3)\) 的逆元。
\end{exercise}
\begin{solution}
若 \(bc-1\in(a)\)，则存在 \(t\in R\) 使 \(1=bc+at\)。任何公因子都整除 \(1\)，所以 \(a,b\) 互素。反过来，\cref{b1:prop:pid-bezout}使互素条件给出 \(1=ua+vb\)，故 \(v+(a)\) 是 \(b+(a)\) 的逆元。

前题已证明 \(3\) 是 \(\mathbb Z[i]\) 中的素元，所以 \((3)\) 是非零素理想。又由\cref{b1:prop:gaussian-euclidean}和\cref{b1:thm:euclidean-pid}，\(\mathbb Z[i]\) 是主理想整环；\cref{b1:prop:pid-irreducible-prime}于是保证 \((3)\) 为极大理想，商环为域。这里 \(N(3)=9\)，可见在 Gauss 整数中，范数为整数素数只是证明不可约的充分条件，不是必要条件。每个剩余类唯一以 \(u+vi\)、\(u,v\in\{0,1,2\}\) 表示，因为相差 \(3\) 的倍数等价于实部与虚部系数分别相差 \(3\) 的倍数，故商有九个元素。最后
\((1+i)(-1+i)=-2\equiv1\pmod3\)，所求逆元为 \(-1+i+(3)\)。
\end{solution}

\begin{exercise}\label{b1:ex:11-squarefree-quotient}
设 \(R\) 为主理想整环，非零非单位元 \(a\) 的不可约分解为 \(a=u p_1^{e_1}\cdots p_t^{e_t}\)，其中 \(u\in R^\times\)，各 \(p_j\) 为两两不相伴的不可约元，\(e_j\geq1\)。证明以下条件等价：商环 \(R/(a)\) 中，每当 \(z^n=0\)、\(n\geq1\)，都有 \(z=0\)；所有指数 \(e_j\) 都等于 \(1\)。
\end{exercise}
\begin{solution}
若全部 \(e_j=1\)，且 \(b^n\in(a)\)，则每个 \(p_j\mid b^n\)。\cref{b1:prop:pid-irreducible-prime}给出 \(p_j\) 的素性，反复用于 \(b^n=b\cdot b^{n-1}\)，得到 \(p_j\mid b\)。唯一分解使各个不同相伴类的因子同时出现于 \(b\)，于是 \(a\mid b\)，即 \(z=b+(a)=0\)。

若某个 \(e_j\geq2\)，取
\[
b=\prod_{j=1}^{t}p_j^{\lceil e_j/2\rceil}.
\]
每个指数在平方后至少达到 \(e_j\)，故 \(a\mid b^2\)。至少一个指数 \(\lceil e_j/2\rceil\) 严格小于 \(e_j\)，所以由分解唯一性，\(a\nmid b\)。因此非零剩余类 \(b+(a)\) 的平方为零，给出反例。
\end{solution}

\begin{exercise}\label{b1:ex:11-no-gcd}
在 \(R=\mathbb Z[\sqrt{-5}]\) 中证明 \(6\) 与 \(2(1+\sqrt{-5})\) 没有最大公因子。
\end{exercise}
\begin{solution}
记 \(s=\sqrt{-5}\)。\(2\) 是这两个元素的公因子；\(1+s\) 也是公因子，因为 \(6=(1+s)(1-s)\)。若存在最大公因子 \(d\)，它首先必须被每个公因子整除，故 \(2\mid d\)、\(1+s\mid d\)。范数的乘法性于是给出
\[
4\mid N(d),\qquad6\mid N(d),\qquad\text{因而 }12\mid N(d).
\]
另一方面，\(d\) 本身必须是公因子，即 \(d\mid6\)、\(d\mid2(1+s)\)，所以
\[
N(d)\mid36,\qquad N(d)\mid24,\qquad\text{因而 }N(d)\mid12.
\]
合并两个方向，正整数 \(N(d)\) 只能等于 \(12\)。但方程 \(a^2+5b^2=12\) 没有整数解：\(\lvert b\rvert\geq2\) 时左端至少为 \(20\)；\(b=0\) 要求 \(a^2=12\)；\(b=\pm1\) 要求 \(a^2=7\)。矛盾，因此不存在最大公因子。
\end{solution}

\begin{exercise}\label{b1:ex:11-coprime-powers}
设 \(R\) 为唯一分解整环，非零元素 \(a,b\) 互素，且 \(ab=c^n\)，其中 \(n\geq2\)。证明存在 \(d,e\in R\) 与单位 \(u,v\)，使 \(a=ud^n\)、\(b=ve^n\)。说明为什么一般不能删去单位因子。
\end{exercise}
\begin{solution}
互素意味着 \(a,b\) 没有公共不可约因子；否则该因子作为公因子应整除 \(1\)，不可能。由 \(ab=c^n\) 及唯一分解，乘积中每个不可约因子的指数都为 \(n\) 的倍数。一个因子只能出现在 \(a,b\) 的一边，所以它在该边的指数本身就是 \(n\) 的倍数。把这些指数分别除以 \(n\)，得到 \(d,e\)；原分解的单位分别保留为 \(u,v\)，即得结论。若某一元素为单位，该边取相应的空乘积 \(d=1\) 或 \(e=1\)。

单位不能普遍删去。例如在 \(\mathbb Z\) 中取 \(a=-1\)、\(b=-4\)、\(c=2\)、\(n=2\)。两者互素且 \(ab=4=c^2\)，但负整数不是整数的平方；乘上单位 \(-1\) 后才得到上述形式。
\end{solution}

\begin{exercise}\label{b1:ex:11-cusp-subring}
设 \(F\) 为域，令
\(A=F[t^2,t^3]\subseteq F[t]\)，即由 \(F\) 中的常数及 \(t^2,t^3\) 通过有限次加法、减法和乘法生成的子环。证明 \(A\) 不是唯一分解整环，尽管 \(F[t]\) 是。
\end{exercise}
\begin{solution}
由于每个整数 \(m\geq2\) 都可写成 \(2r+3s\)、\(r,s\geq0\)，环 \(A\) 恰由一次项系数为零的多项式组成。它的单位是非零常数：若两个多项式乘积为 \(1\)，次数相加迫使两者都是常数。于是每个非零非单位的次数至少为 \(2\)。因此 \(t^2,t^3\) 都不可约，否则两个非单位因子的乘积次数至少为 \(4\)。它们不相伴，因为乘以非零常数不改变次数。但
\[
t^6=(t^2)^3=(t^3)^2
\]
给出长度不同的两种不可约分解，所以 \(A\) 不是唯一分解整环。另一方面，\cref{b1:prop:field-poly-division}给出 \(F[t]\) 为 Euclid 整环，再由\cref{b1:thm:euclidean-pid}与\cref{b1:thm:pid-ufd}得它是唯一分解整环。同一个元素 \(t^3\) 在 \(A\) 中不可约，在 \(F[t]\) 中却可分解为 \(t\cdot t^2\)，因为 \(t\notin A\)。可用的因子随所在环改变，唯一分解性不能任意传给子环。
\end{solution}

\begin{exercise}\label{b1:ex:11-pid-quotient-ideals}
设 \(R\) 为主理想整环，\(a=u p_1^{e_1}\cdots p_t^{e_t}\) 为非零非单位的不可约分解，其中各 \(p_j\) 两两不相伴、\(e_j\geq1\)。描述 \(R/(a)\) 的全部理想及包含关系，并求理想总数与极大理想总数。
\end{exercise}
\begin{solution}
由\cref{b1:prop:ring-ideal-correspondence}，商环的理想与 \(R\) 中包含 \((a)\) 的理想一一对应。后者均为 \((d)\)，且由\cref{b1:prop:associate-ideals}，\((a)\subseteq(d)\) 等价于 \(d\mid a\)。\cref{b1:thm:pid-ufd}的唯一分解使每个这样的 \(d\) 在相伴意义下唯一写为
\[
d=\prod_{j=1}^t p_j^{f_j},\qquad 0\leq f_j\leq e_j.
\]
不同的指数列给出不同的主理想，因为主理想相等恰好等价于生成元相伴。因此商环全部理想为这些 \((d)/(a)\)，总数为 \(\prod_{j=1}^t(e_j+1)\)。

若 \(d'\) 的指数列为 \(f'_j\)，则
\((d)/(a)\subseteq(d')/(a)\) 当且仅当 \(d'\mid d\)，也就是对每个 \(j\) 都有 \(f'_j\leq f_j\)。全环对应全零指数列，极大真理想恰对应只在一个位置取 \(1\)、其余取 \(0\) 的指数列。它们即 \((p_j)/(a)\)，共有 \(t\) 个；若指数和大于 \(1\)，减去一个正指数就会得到严格位于该理想与全环之间的理想，因而不极大。

包含关系与指数大小的方向相反。例如 \(a=p^3\)、\(p\) 不可约时，商环中的全部理想组成严格升链
\[
(p^3)/(p^3)\subsetneq(p^2)/(p^3)
\subsetneq(p)/(p^3)\subsetneq R/(p^3).
\]
从左到右，生成元的指数依次为 \(3,2,1,0\)，而理想依次变大。
\end{solution}

\begin{optionalexercise}\label{b1:ex:11-nineteen-quotient}
沿用\subsecref{b1:subsec:1103:03}中的 \(R=\mathbb Z[\omega]\)、\(\omega=(1+\sqrt{-19})/2\)，取 \(N(z)=\lvert z\rvert^2\)。求商环 \(R/(2)\)，证明它为四元素域，并求 \(\omega+(2)\) 的逆元。另证不存在 \(q,r\in R\) 使 \(\omega=2q+r\)、\(N(r)<N(2)\)。说明这项除法失败本身为何还不能排除别的 Euclid 函数。
\end{optionalexercise}
\begin{solution}
因为每个元素唯一写为 \(a+b\omega\)，模 \((2)\) 后有四个不同的类
\(0,1,\overline\omega,1+\overline\omega\)。关系 \(\omega^2=\omega-5\) 在商中成为
\(\overline\omega^2=\overline\omega+1\)，所以
\[
\overline\omega(1+\overline\omega)=1.
\]
三个非零类 \(1,\overline\omega,1+\overline\omega\) 的逆元分别为 \(1,1+\overline\omega,\overline\omega\)。故商是四元素域；所求逆元为 \(1+\omega+(2)\)。这也显式给出同构
\(R/(2)\cong\mathbb F_2[t]/(t^2+t+1)\)：把 \(t\) 的类送到 \(\overline\omega\)，两边均以 \(1,t\) 或 \(1,\overline\omega\) 唯一表示元素，关系相同。

若 \(q=a+b\omega\)，则 \(r=\omega-2q=-2a+(1-2b)\omega\)。记 \(v=1-2b\)，它是非零奇数。由 \(\omega=1/2+(\sqrt{19}/2)i\) 直接计算实部与虚部，得
\[
N(r)=\Bigl(-2a+\frac v2\Bigr)^2+\frac{19v^2}{4}
\geq\frac{19}{4}>4=N(2).
\]
所以没有所要求的范数余数。这个计算只否定函数 \(N\) 的 Euclid 条件；环的 Euclid 性要求存在某个合适函数，不能据此否定全部可能性。排除任意 Euclid 函数的结论另由\cref{b1:prop:noneuclidean-pid}给出，其证明使用了最小 Euclid 值的非单位元以及商环大小的限制。
\end{solution}

\begin{exercise}\label{b1:ex:11-nonprincipal-square}
在 \(\mathbb Z[x]\) 中令 \(J=(2,x)\)。对每个整数 \(n\geq1\)，求 \(J^n\) 的一组生成元，证明 \(J^n\) 不是主理想，并求 \(\sqrt{J^n}\)。说明为何发现非主理想不能据此断言该环不是唯一分解整环。
\end{exercise}
\begin{solution}
由理想积的有限和定义，归纳可得
\[
J^n=(2^n,2^{n-1}x,\ldots,2x^{n-1},x^n).
\]
确实，右边在 \(n=1\) 时就是 \(J\)；乘以 \((2,x)\) 后，每个生成元 \(2^{n-j}x^j\) 分别给出 \(2^{n+1-j}x^j\) 与 \(2^{n-j}x^{j+1}\)，恰好覆盖下一次幂的全部生成元。这些生成元及其任意多项式系数组合的常数项都被 \(2^n\) 整除，故 \(1\notin J^n\)。若 \(J^n=(h)\)，则 \(h\mid2^n\)。由\cref{b1:prop:polynomial-degree}，\(h\) 为非零整数常数；再由 \(h\mid x^n\)，比较 \(x^n\) 的系数，得 \(h\mid1\)，所以 \(h=\pm1\)。这将给出 \(J^n=\mathbb Z[x]\)，矛盾。

映射 \(\mathbb Z[x]\rightarrow\mathbb F_2\)、\(f\mapsto\overline{f(0)}\) 是满环同态。它的核恰为常数项为偶数的多项式集合，也就是 \(J=(2,x)\)。由\cref{b1:thm:ring-first-iso}和\cref{b1:prop:prime-max-quotient}，\(J\) 为极大理想，因而为素理想。若 \(f^m\in J\)，反复应用素性就得 \(f\in J\)，故 \(\sqrt J=J\)。再由\cref{b1:prop:ideal-containments}中 \(\sqrt{J^n}=\sqrt J\) 的等式，得到 \(\sqrt{J^n}=J\)。

理想 \(J^n\) 的非主性否定的是主理想整环条件。唯一分解整环的定义要求元素的不可约分解存在且唯一，并不要求全部理想为主理想；因此本题的非主性证明不能否定它。下一章\cref{b1:thm:polynomial-ufd}将说明 \(\mathbb Z[x]\) 实际上是唯一分解整环。
\end{solution}

\begin{exercise}\label{b1:ex:11-factorization-nonexistence}
设 \(F\) 为域。令 \(A\) 由形式有限和 \(\sum_q a_qT^q\) 组成，其中 \(a_q\in F\)，指数 \(q\) 取形如 \(m/2^n\) 的非负数，\(m,n\in\mathbb Z_{\geq0}\)。两个形式和相等指每个指数的系数相等；加法逐项进行，乘法由 \(T^qT^r=T^{q+r}\) 及分配律规定。证明 \(A\) 是整环，且 \(T\) 不能分解为有限个不可约元的乘积。再证明 \(A\) 的每个有限生成理想都是主理想，\(A\) 的每个不可约元都是素元，并说明为何 \(A\) 仍不是主理想整环。由此区分分解不存在与分解不唯一两种失败。
\end{exercise}
\begin{solution}
本题要分别检验两种容易被误用的推论：每个不可约元都是素元，仍不能代替不可约分解的存在性；每个有限生成理想都是主理想，也不能代替所有理想都是主理想。下面先考察元素的分解，再考察理想。

任意有限多个这样的形式和，其涉及的指数都可取共同分母 \(2^n\)；以 \(U=T^{1/2^n}\) 表示时，它们就是通常的 \(F[U]\) 中的多项式。运算结果仍属 \(A\)，加法、乘法的环公理可在这个共同多项式环中验证。因此 \(A\) 是交换含幺环。

对非零 \(f\in A\)，记其最大与最小非零系数指数分别为 \(d(f),v(f)\)。因 \(F\) 为域，最高和最低指数处的系数相乘均非零，故对非零 \(f,g\)，有
\[
d(fg)=d(f)+d(g),\qquad v(fg)=v(f)+v(g).
\]
这首先说明 \(fg\ne0\)，所以 \(A\) 为整环；又说明其单位恰为非零常数，因为 \(fg=1\) 迫使两个非负的最大指数都为零。

若 \(T=fg\)，则
\[
0=d(T)-v(T)
=\bigl(d(f)-v(f)\bigr)+\bigl(d(g)-v(g)\bigr).
\]
两个括号内均为非负数，故各自为零，\(f,g\) 都是单项式。因此 \(T\) 的每个非单位因子都是 \(cT^q\)、\(c\ne0,q>0\)。但
\[
cT^q=(cT^{q/2})T^{q/2}
\]
总是两个非单位的乘积，因为 \(q/2\) 仍是允许的正指数。所以 \(T\) 没有不可约因子，更不可能有有限不可约分解。

有限多个元素 \(f_1,\ldots,f_r\in A\) 同属于某个 \(F[U]\)，其中 \(U=T^{1/2^n}\)。由于 \(F[U]\) 是主理想整环，存在 \(d\in F[U]\) 使 \((f_1,\ldots,f_r)F[U]=dF[U]\)。于是每个 \(f_j\) 都是 \(d\) 在 \(A\) 中的倍数，而 \(d\) 是 \(f_1,\ldots,f_r\) 在 \(F[U]\) 中的线性组合，也属于它们在 \(A\) 中生成的理想。因此 \((f_1,\ldots,f_r)A=dA\)；零理想也由 \(0\) 生成。

设 \(p\in A\) 不可约，且 \(p\mid ab\)。若 \(p\mid a\)，素性所需的结论已经成立；现设 \(p\nmid a\)。由于二生成理想为主理想，可写 \((p,a)=(d)\)。由 \(d\mid p\) 和 \(p\) 的不可约性，\(d\) 或为单位，或与 \(p\) 相伴；后一种情形会使 \(p\mid a\)，故 \((p,a)=A\)。于是存在 \(u,v\in A\)，使 \(1=up+va\)。两边乘 \(b\)，再用 \(p\mid ab\)，即得 \(p\mid b\)。所以 \(p\) 是素元。

现在把不断分解元素与不断增大主理想并排来看：
\[
T=(T^{1/2})^2=(T^{1/4})^4=\cdots,
\qquad
(T)\subsetneq(T^{1/2})\subsetneq(T^{1/4})\subsetneq\cdots.
\]
每个包含都严格，因为 \((T^{1/2^n})\) 的非零元素的最小指数至少为 \(1/2^n\)，而 \(T^{1/2^{n+1}}\) 的最小指数更小。\cref{b1:thm:pid-ufd}证明分解存在性时所用的主理想升链稳定性，在这里失效了。还可直接构造一个非主理想：令
\[
J=\bigcup_{n\geq0}(T^{1/2^n}).
\]
这是理想升链的并，因而是理想。若 \(J\) 有有限个生成元，它们必同时落在某个 \((T^{1/2^N})\) 中，便有 \(J\subseteq(T^{1/2^N})\)；但下一项的生成元属于 \(J\)，却不属于该理想，矛盾。因此 \(J\) 不是有限生成理想，尤其不是主理想，\(A\) 不是主理想整环。这也具体说明了为何前面关于有限生成理想的结论仍然成立。

即使每个不可约元都是素元，也不能省略唯一分解整环所要求的分解存在性。与\cref{b1:prop:sqrt-minus-five-factorization}中 \(6\) 有两种不等价不可约分解相比，本例失败的是分解存在性。
\end{solution}

\begin{exercise}\label{b1:ex:11-ideal-splitting-three}
令 \(s=\sqrt{-5}\)、\(R=\mathbb Z[s]\subseteq\mathbb C\)，考虑映射
\[
\Phi:R\longrightarrow\mathbb F_7\times\mathbb F_7,\qquad
a+bs\longmapsto(\overline{a+3b},\overline{a-3b}).
\]
证明 \(\Phi\) 是保持单位的满环同态，且 \(\ker\Phi=(7)\)。按第一、第二坐标的顺序，将两个坐标映射的核记为 \(P,Q\)。求它们的理想生成元，并证明它们是不同的极大理想，且 \(P+Q=R\)、\(P\cap Q=PQ=(7)\)。确定 \(R/(7)\) 的全部幂等元，说明两个非平凡幂等元分别怎样选取上述直积中的一个分量。最后证明 \(7\) 不可约却不是素元，而 \(P,Q\) 都不是主理想。
\end{exercise}
\begin{solution}
两种模 \(7\) 的求值 \(s\mapsto3\)、\(s\mapsto-3\) 都保持关系 \(s^2=-5\)，因为 \((\pm3)^2=9\equiv-5\pmod7\)。所以 \(\Phi\) 保持加法、乘法和单位。任给 \(u,v\in\mathbb F_7\)，选整数 \(a,b\) 满足
\[
\bar a=4(u+v),\qquad \bar b=6(u-v).
\]
这里 \(4=2^{-1}\)、\(6=6^{-1}\) 在 \(\mathbb F_7\) 中成立，代入得到 \(\Phi(a+bs)=(u,v)\)，证明满射。若 \(\Phi(a+bs)=0\)，两坐标相加、相减给出 \(2\bar a=0\)、\(6\bar b=0\)，故 \(7\mid a\)、\(7\mid b\)。反过来，\(7\) 的倍数都在核中，因此 \(\ker\Phi=(7)\)。由\cref{b1:thm:ring-first-iso}，\(\Phi\) 诱导同构
\[
\bar\Phi:R/(7)\longrightarrow\mathbb F_7\times\mathbb F_7.
\]

第一坐标为零的条件是 \(7\mid a+3b\)，此时 \(a+bs=b(s-3)+(a+3b)\)，所以 \(P=(7,s-3)\)；同理，由 \(a+bs=b(s+3)+(a-3b)\) 得 \(Q=(7,s+3)\)。两个坐标映射都满射到域，由\cref{b1:thm:ring-first-iso}分别得到 \(R/P\cong\mathbb F_7\)、\(R/Q\cong\mathbb F_7\)，再由\cref{b1:prop:prime-max-quotient}知 \(P,Q\) 都是极大理想。元素 \(s-3\) 属于 \(P\)，其第二坐标为 \(-6=1\ne0\)，故 \(P\ne Q\)。又有 \(6=(s+3)-(s-3)\in P+Q\) 及 \(7\in P+Q\)，所以 \(1=7-6\in P+Q\)。由\cref{b1:lem:comaximal-product}，\(PQ=P\cap Q\)；而两坐标核的交就是 \(\ker\Phi\)，故 \(PQ=P\cap Q=(7)\)。

域中的幂等元只有 \(0,1\)，所以直积中的幂等元恰为 \((0,0),(1,1),(1,0),(0,1)\)。它们在 \(R/(7)\) 中的原像分别为
\[
0,\qquad1,\qquad e_+=(4-s)+(7),\qquad e_-=(s-3)+(7).
\]
确实，\(\bar\Phi(e_+)=(1,0)\)、\(\bar\Phi(e_-)=(0,1)\)。若 \(\bar\Phi(z)=(u,v)\)，则 \(\bar\Phi(e_+z)=(u,0)\)、\(\bar\Phi(e_-z)=(0,v)\)：乘以这两个幂等元分别保留第一、第二分量。因而 \(e_++e_-=1\)、\(e_+e_-=0\)，且它们生成的两个理想分别对应直积的两个分量。

\(N(7)=49\)。若 \(7\) 有非平凡分解，两个非单位因子的范数只能都为 \(7\)。但 \(a^2+5b^2=7\) 无整数解：\(b=0\) 要求 \(a^2=7\)，\(|b|=1\) 要求 \(a^2=2\)，\(|b|\ge2\) 则左边至少为 \(20\)。所以 \(7\) 不可约。另一方面，
\[
(s-3)(s+3)=s^2-9=-14
\]
被 \(7\) 整除，两个因子却都不被 \(7\) 整除，因为它们的 \(s\) 系数为 \(1\)。故 \(7\) 不是素元。

若 \(P=(d)\)，由 \(7\in P\) 得 \(d\mid7\)。不可约性说明 \(d\) 或为单位，或与 \(7\) 相伴。前者与 \(P\) 为真理想矛盾；后者使 \(P=(7)\)，与 \(s-3\in P\setminus(7)\) 矛盾。因此 \(P\) 非主，对 \(Q\) 用 \(s+3\) 作同样判断即可。这样，\((7)=PQ\) 是非主理想之间的乘积分解，与元素 \(7\) 的不可约性相容。
\end{solution}
